\begin{abstract}
Line heating is a practical shipyard forming operation for producing curvature in steel hull plates, but process planning still relies heavily on operator experience because residual shape is difficult to predict quantitatively. This paper presents a reproducible three-dimensional thermo-mechanical finite element workflow for ship-plate line heating. The workflow converts a prescribed heating plan into transient temperature fields and post-cooling deformation using temperature-dependent Q345 material properties, moving Gaussian and full-line thermal-flux source options, convection, quenching, optional radiation, and a linear thermoelastic mechanical solve augmented by a calibrated inherent-strain surrogate for residual bending. The reported residual-shape validation is therefore a validation of this reduced residual-deformation workflow, not a claim of fully resolved transient elastoplastic residual-stress prediction.

The workflow is assessed against an available subset of eleven published Q345 line-heating benchmark cases. The inherent-strain parameters are calibrated on four anchor cases and then held fixed for the remaining non-anchor cases. For the moving-torch baseline, the mean absolute percentage error in midspan residual deflection is \SI{1.17}{\%} over all cases and \SI{1.55}{\%} on the non-anchor validation cases. A process-scheduling study shows that replacing two separate passes by one collapsed full-line heating event underpredicts multi-pass deflection by approximately \SI{50}{\%}; preserving the sequential thermal cycles recovers the multi-pass response with MAPE \SI{0.27}{\%} on the evaluated subset. A keel-plate numerical demonstration illustrates how the same workflow can be used for application-scale ship-production planning, but it is not treated as experimental validation. The results show that accurate representation of pass scheduling is as important as total heat input for multi-pass line-heating prediction.
\end{abstract}
